% ============================================================================
% CAPÍTULO 7: DISCUSIÓN
% ============================================================================

\chapter{Discusión}
\label{cap:discusion}

\justifying
\setlength{\parskip}{0pt}

% ----------------------------------------------------------------------------
% INSTRUCCIONES PARA ESCRIBIR LA DISCUSIÓN:
% ----------------------------------------------------------------------------
% - INTERPRETA los resultados (no solo los repitas)
% - COMPARA con estudios previos (de tus antecedentes)
% - EXPLICA similitudes y diferencias
% - Identifica LOGROS y LIMITACIONES
% - Proporciona EXPLICACIONES teóricas
% - Extensión típica: 15-20 páginas
% ----------------------------------------------------------------------------

{[Párrafo introductorio que sintetiza los hallazgos principales]}

En este capítulo se interpretan y discuten los resultados presentados en el capítulo anterior, contrastándolos con la evidencia científica disponible y el marco teórico establecido.

% ============================================================================
% SECCIONES SUGERIDAS PARA ESTRUCTURAR LA DISCUSIÓN
% ============================================================================

% ----------------------------------------------------------------------------
% SECCIÓN 7.1: Discusión de Hallazgos Principales
% ----------------------------------------------------------------------------
\section{Discusión de Hallazgos Principales}
\label{sec:discusion_hallazgos}

\subsection{{[Hallazgo Principal 1]}}
\label{subsec:discusion_hallazgo1}

% ESTRUCTURA SUGERIDA PARA CADA HALLAZGO:
% 1. Recuerda brevemente el hallazgo
% 2. Interpreta: ¿qué significa?
% 3. Compara con la literatura
% 4. Explica similitudes/diferencias
% 5. Proporciona explicaciones teóricas

\textbf{Hallazgo:}

En este estudio se encontró que \textit{{[resume brevemente el hallazgo principal 1]}}.

\textbf{Interpretación:}

Este hallazgo sugiere que \textit{{[tu interpretación del significado]}}, lo cual implica \textit{{[implicaciones]}}.

\textbf{Comparación con la literatura:}

Este resultado es \textit{{[consistente/inconsistente]}} con lo reportado por \cite{autor1}, quien encontró que \textit{{[qué encontró]}}. Similarmente, \cite{autor2} reportó \textit{{[hallazgo similar/diferente]}}.

Sin embargo, difiere de los hallazgos de \cite{autor3}, quienes reportaron \textit{{[diferencia]}}. Esta discrepancia podría explicarse por \textit{{[posibles razones: diferencias en población, metodología, contexto, etc.]}}.

\textbf{Explicación teórica:}

Desde la perspectiva de \textit{{[teoría o modelo del marco teórico]}}, este resultado se explica porque \textit{{[explicación basada en teoría]}} \cite{referencia_teoria}.

\subsection{{[Hallazgo Principal 2]}}
\label{subsec:discusion_hallazgo2}

{[Repite la estructura para cada hallazgo principal]}

\subsection{{[Hallazgo Principal 3]}}
\label{subsec:discusion_hallazgo3}

{[Continúa con todos tus hallazgos principales]}

% ----------------------------------------------------------------------------
% SECCIÓN 7.2: Integración de Resultados con el Cuadro Comparativo
% ----------------------------------------------------------------------------
\section{Comparación con Investigaciones Previas}
\label{sec:comparacion_investigaciones}

% INSTRUCCIONES:
% - Retoma el cuadro comparativo del Cap. 2 (Marco Teórico)
% - Discute cómo tus resultados se comparan con esos estudios

\subsection{Síntesis Comparativa}
\label{subsec:sintesis_comparativa}

La \Cref{tab:sintesis_discusion} presenta una síntesis comparativa entre los hallazgos de estudios previos y los resultados de esta investigación.

\begin{table}[htbp]
    \centering
    \caption{Comparación de hallazgos principales con investigaciones previas}
    \label{tab:sintesis_discusion}
    \small
    \begin{tabular}{@{}p{3cm}p{4cm}p{4cm}p{3cm}@{}}
        \toprule
        \textbf{Estudio} & \textbf{Hallazgo Principal} & \textbf{Hallazgo de Esta Investigación} & \textbf{Consistencia} \\
        \midrule
        
        Autor1 et al. (2023) & 
        {[Hallazgo del estudio]} & 
        {[Tu hallazgo relacionado]} & 
        \textit{Consistente/ Parcialmente/ Inconsistente} \\
        \midrule
        
        Autor2 \& Autor3 (2022) & 
        {[Hallazgo]} & 
        {[Tu hallazgo]} & 
        \textit{{[Consistencia]}} \\
        \midrule
        
        {[Continúa con más estudios]} & 
        ... & 
        ... & 
        ... \\
        
        \bottomrule
    \end{tabular}
\end{table}

\subsection{Análisis de Convergencias}
\label{subsec:convergencias}

Los resultados de esta investigación convergen con la literatura previa en los siguientes aspectos:

\begin{enumerate}
    \item \textbf{Aspecto 1:} Al igual que \cite{autor1} y \cite{autor2}, este estudio confirma que \textit{{[punto de convergencia]}}. Esto sugiere que \textit{{[implicación]}}.
    
    \item \textbf{Aspecto 2:} \textit{{[Otra convergencia]}}...
    
    \item \textbf{Aspecto 3:} \textit{{[Otra convergencia]}}...
\end{enumerate}

\subsection{Análisis de Divergencias}
\label{subsec:divergencias}

Por otro lado, se identificaron diferencias con algunos estudios previos:

\begin{enumerate}
    \item \textbf{Diferencia 1:} A diferencia de \cite{autor3}, quien reportó \textit{{[hallazgo diferente]}}, en este estudio se encontró \textit{{[tu hallazgo]}}. 
    
    \textbf{Posibles explicaciones:}
    \begin{itemize}
        \item Diferencias en la población estudiada: \textit{{[explica]}}
        \item Diferencias metodológicas: \textit{{[explica]}}
        \item Diferencias contextuales: \textit{{[explica]}}
    \end{itemize}
    
    \item \textbf{Diferencia 2:} \textit{{[Otra divergencia con explicación]}}...
\end{enumerate}

% ----------------------------------------------------------------------------
% SECCIÓN 7.3: Logros de la Investigación
% ----------------------------------------------------------------------------
\section{Logros de la Investigación}
\label{sec:logros}

Esta investigación logró:

\begin{enumerate}
    \item \textbf{Logro 1:} Se cumplió el objetivo de \textit{{[objetivo]}} al \textit{{[cómo se logró]}}, lo cual representa una aportación importante porque \textit{{[por qué es importante]}}.
    
    \item \textbf{Logro 2:} Se confirmó la hipótesis de que \textit{{[hipótesis]}}, con evidencia estadística significativa ($p < $ \textit{{[valor]}}), lo cual \textit{{[implicación]}}.
    
    \item \textbf{Logro 3:} Se logró \textit{{[validar/desarrollar/adaptar]}} el instrumento \textit{{[nombre]}} en el contexto de \textit{{[población]}}, con adecuadas propiedades psicométricas ($\alpha$ = \textit{{[valor]}}), lo cual contribuye metodológicamente al campo.
    
    \item \textbf{Logro 4:} Se identificaron predictores significativos de \textit{{[variable]}}, específicamente \textit{{[predictores]}}, explicando el \textit{[\%]} de la varianza, lo cual tiene implicaciones prácticas para \textit{{[aplicación]}}.
    
    \item \textbf{Logro 5:} Se generó conocimiento nuevo sobre \textit{{[aspecto específico]}} que no había sido explorado previamente en \textit{{[contexto/población]}}.
\end{enumerate}

% ----------------------------------------------------------------------------
% SECCIÓN 7.4: Limitaciones del Estudio
% ----------------------------------------------------------------------------
\section{Limitaciones del Estudio}
\label{sec:limitaciones}

% INSTRUCCIONES:
% - Sé HONESTO sobre las limitaciones
% - Esto demuestra rigor científico, NO debilita tu trabajo
% - Categoriza las limitaciones (metodológicas, contextuales, muestrales, etc.)

A pesar de los logros alcanzados, es importante reconocer las siguientes limitaciones:

\subsection{Limitaciones Metodológicas}
\label{subsec:limitaciones_metodologicas}

\begin{enumerate}
    \item \textbf{Diseño del estudio:} El diseño \textit{{[transversal/observacional/etc.]}} no permite establecer relaciones causales, únicamente asociaciones. Para determinar causalidad se requeriría un diseño \textit{{[experimental/longitudinal]}}.
    
    \item \textbf{Instrumentos:} La recolección de datos se basó en \textit{{[autorreporte/cuestionarios]}}, lo cual puede estar sujeto a sesgos de \textit{{[deseabilidad social/memoria/etc.]}}. Mediciones objetivas (como \textit{{[ejemplos]}}) podrían complementar los hallazgos.
    
    \item \textbf{Variables confusoras:} Aunque se controlaron \textit{{[variables controladas]}}, es posible que otras variables no medidas (como \textit{{[ejemplos]}}) puedan estar influyendo en los resultados.
\end{enumerate}

\subsection{Limitaciones Muestrales}
\label{subsec:limitaciones_muestrales}

\begin{enumerate}
    \item \textbf{Tamaño de muestra:} Si bien la muestra fue suficiente para detectar efectos significativos, un tamaño muestral mayor ($n > $ \textit{{[número]}}) podría permitir análisis de subgrupos y aumentar el poder estadístico.
    
    \item \textbf{Tipo de muestreo:} El muestreo \textit{{[tipo]}} puede limitar la generalización de los resultados a \textit{{[población más amplia]}}. Estudios futuros con muestreo probabilístico fortalecerían la validez externa.
    
    \item \textbf{Características de la muestra:} La muestra estuvo compuesta principalmente por \textit{{[característica, ej: adultos jóvenes, mujeres, etc.]}}, lo cual puede limitar la aplicabilidad a \textit{{[otros grupos]}}.
\end{enumerate}

\subsection{Limitaciones Contextuales}
\label{subsec:limitaciones_contextuales}

\begin{enumerate}
    \item \textbf{Contexto geográfico:} Los resultados provienen de \textit{{[Chihuahua/contexto específico]}}, por lo que la generalización a otras regiones con características \textit{{[sociodemográficas/culturales/etc.]}} diferentes debe hacerse con cautela.
    
    \item \textbf{Período temporal:} Los datos se recolectaron durante \textit{{[período]}}, lo cual podría haber sido influenciado por \textit{{[eventos contextuales como pandemia, estacionalidad, etc.]}}.
\end{enumerate}

\subsection{Limitaciones de Recursos}
\label{subsec:limitaciones_recursos}

\begin{enumerate}
    \item \textbf{Recursos económicos:} Los recursos disponibles limitaron \textit{{[aspecto, ej: tamaño de muestra, número de mediciones, uso de tecnología específica]}}.
    
    \item \textbf{Tiempo:} El tiempo disponible para el estudio (\textit{{[duración]}}) no permitió \textit{{[seguimiento longitudinal/mediciones repetidas/etc.]}}.
\end{enumerate}

% ----------------------------------------------------------------------------
% SECCIÓN 7.5: Riesgos Identificados
% ----------------------------------------------------------------------------
\section{Riesgos Identificados}
\label{sec:riesgos}

Durante el desarrollo de esta investigación se identificaron los siguientes riesgos:

\begin{enumerate}
    \item \textbf{Sesgo de selección:} \textit{{[Describe el riesgo y cómo se mitigó]}}
    
    \item \textbf{Sesgo de respuesta:} \textit{{[Describe el riesgo y cómo se mitigó]}}
    
    \item \textbf{Pérdida de participantes:} \textit{{[Si aplica, describe la tasa de pérdida y su posible impacto]}}
\end{enumerate}

% ----------------------------------------------------------------------------
% SECCIÓN 7.6: Implicaciones de los Hallazgos
% ----------------------------------------------------------------------------
\section{Implicaciones de los Hallazgos}
\label{sec:implicaciones}

\subsection{Implicaciones Teóricas}
\label{subsec:implicaciones_teoricas}

Los resultados de esta investigación tienen las siguientes implicaciones teóricas:

\begin{enumerate}
    \item Apoyan la teoría de \textit{{[teoría]}} al confirmar que \textit{{[aspecto]}}
    \item Sugieren una revisión de \textit{{[concepto/modelo]}} en el sentido de \textit{{[sugerencia]}}
    \item Aportan evidencia sobre \textit{{[mecanismo/proceso]}} que no había sido claramente demostrado
\end{enumerate}

\subsection{Implicaciones Prácticas}
\label{subsec:implicaciones_practicas}

Los hallazgos tienen aplicaciones prácticas inmediatas:

\begin{enumerate}
    \item \textbf{Para profesionales de la salud:} Los resultados sugieren que \textit{{[recomendación práctica específica]}}
    
    \item \textbf{Para tomadores de decisiones:} Se recomienda \textit{{[política/programa]}} basándose en \textit{{[hallazgo]}}
    
    \item \textbf{Para la población objetivo:} \textit{{[Implicación para beneficiarios]}}
\end{enumerate}

\subsection{Implicaciones Metodológicas}
\label{subsec:implicaciones_metodologicas}

Metodológicamente, este estudio aporta:

\begin{enumerate}
    \item Un instrumento validado para \textit{{[propósito]}} en \textit{{[población]}}
    \item Una estrategia metodológica para \textit{{[proceso]}}
    \item Evidencia sobre la utilidad de \textit{{[técnica/enfoque]}}
\end{enumerate}

% ----------------------------------------------------------------------------
% SECCIÓN 7.7: Reflexión sobre el Cumplimiento de Objetivos e Hipótesis
% ----------------------------------------------------------------------------
\section{Cumplimiento de Objetivos e Hipótesis}
\label{sec:cumplimiento}

\subsection{Objetivo General}
\label{subsec:cumplimiento_og}

El objetivo general de \textit{{[enunciado]}} fue alcanzado satisfactoriamente al \textit{{[cómo se logró]}}.

\subsection{Objetivos Específicos}
\label{subsec:cumplimiento_oe}

\begin{itemize}
    \item \textbf{OE1:} \textit{{[Enunciado]}} - \textbf{Cumplido.} \textit{{[Evidencia del cumplimiento]}}
    \item \textbf{OE2:} \textit{{[Enunciado]}} - \textbf{Cumplido.} \textit{{[Evidencia]}}
    \item \textbf{OE3:} \textit{{[Enunciado]}} - \textbf{Cumplido.} \textit{{[Evidencia]}}
    \item \textbf{OE4:} \textit{{[Enunciado]}} - \textbf{Cumplido.} \textit{{[Evidencia]}}
\end{itemize}

\subsection{Hipótesis}
\label{subsec:cumplimiento_hipotesis}

\textbf{Hipótesis de Investigación (Hi):} \textit{{[Enunciado]}}

\textbf{Resultado:} La hipótesis fue \textit{{[confirmada/rechazada]}} con base en \textit{{[evidencia estadística: prueba utilizada, valor estadístico, p-value]}}.

\textit{{[Interpreta qué significa este resultado]}}

% ============================================================================
% NOTAS FINALES:
% ============================================================================
%
% RECUERDA:
% 1. INTERPRETA, no solo describas
% 2. COMPARA con estudios de tus antecedentes
% 3. EXPLICA similitudes y diferencias
% 4. Sé HONESTO con las limitaciones
% 5. Proporciona EXPLICACIONES TEÓRICAS
% 6. Conecta con implicaciones PRÁCTICAS
% 7. Verifica cumplimiento de objetivos e hipótesis
%
% ============================================================================

